\scanpage{155}
Ainsi, $A_1\ast A_2$ est une opération interne dans les espaces affines à involution. Je dis que cette opération a une structure d'associativité, avec \inserted{objet} unité.
% slanted
\margin{il y a une structure de commutativité.}
Ce dernier est l'espace affine trivial $k$ lui-même, avec l'involution $x\longmapsto1-x$, \deleted{\unsure{associée à l'espace}} \inserted{(et non $x\longrightarrow-x$, qui devient triviale en car. $2$ !)}, l'enveloppe affine est $\deleted{\unsure{k\longrightarrow k}}\ k\times k\xrightarrow[\text{somme}]{\varepsilon}k,\quad t_0=(1,1)\mapsto2$.
\inserted{$\sigma_0$ est $(\lambda,\mu)\longmapsto(\mu,\lambda)$}
\margin{$\begin{aligned}k&\hookrightarrow k\times k\\ \lambda&\overset{\mathrm{aff}}{\longmapsto}(\lambda,1-\lambda)\\ k&\overset{\mathrm{lin}}{\hookrightarrow}k\times k\\ \lambda&\longmapsto(\lambda,-\lambda)\end{aligned}$ \inserted{image $A_0=\operatorname{Ker}\unsure{\varepsilon|k}$}}
L'élément $e_0=(\deleted{\unsure{1}}0,1)$ de cet espace affine $A_0$ satisfait
\[
 \sigma_0(e_0)-e_0=1\in k=A_0^0.
\]
Ceci dit, on utilise $e_0$ pour définir une iso. can.
\[
 \begin{aligned}
 A&\overset{\sim}{\longrightarrow} A_0\ast A\deleted{\longrightarrow A}\\
 x&\longmapsto e_0\ast x
 \end{aligned}
\]
compatible avec $A^0\overset{\sim}{\longrightarrow}\underset{\substack{\Vert\\A_0^0}}{k}\otimes A^0$ iso. can.

Pour définir la structure d'associativité, on va définir \unsure{directemt|directement}, pour un nb quelc. de facteurs \inserted{affines involutifs} le produit
\[
 A_1\ast A_2\cdots\ast A_n
\]
% archive sheet number: 154

